% HD-N27. GPT-6 Astra Pro /root/imprimitive, 21 November 2026.
% Inherited effort; numerical setting not exposed.
\section{HD-N27: ample-twisted restriction complexes on the graph union}
\label{sec:ample-twisted-graph-restrictions}

Take any nodal arrangement $Z=\bigcup_{\nu=1}^{9408}C_\nu$
of HD-N23, with embeddings $j_\nu:C_\nu\hookrightarrow A$.
Choose an actual ample line bundle $M_\nu$ on each component
and an actual nonzero section $s_\nu\in H^0(C_\nu,M_\nu)$.
Put
\[
 Q=\bigoplus_\nu j_{\nu,*}M_\nu,\qquad
 q:\mathcal O_A\longrightarrow Q,
       \quad q(1)=(s_\nu)_\nu,\qquad
 F=[\mathcal O_A\xrightarrow{q}Q],                         \tag{N27.1}
\]
where the terms of $F$ are in degrees zero and one. Thus
there is the actual fiber triangle
\[
                   F\xrightarrow{v}\mathcal O_A
                    \xrightarrow{q}Q\longrightarrow F[1].
                                                               \tag{N27.2}
\]
These are perfect gluable simple objects. Every actual
coherent faithful finite point action on any such $F$ has
a weak degree-two trace kernel of dimension at least $13$.
The same bound holds after an arbitrary actual isogeny
pullback, for every coherent point action that then exists,
and after subsequent actual line tensor and Fourier--Mukai
equivalences.

The uniform choices are explicit whole-character candidates.
Let $L$ be the actual polarization line bundle of HD-N23,
$c_1(L)=\theta$, and choose $M_\nu=L^k|_{C_\nu}$, $k>0$.
The nonzero sections exist by the ample abelian vanishing
and Riemann--Roch used below. Projection formula and the
abelian embeddings give the complete character
\[
 \begin{split}
 \operatorname{ch}(F)
 &=1-e^{k\theta}(32\theta^3+120T)\\
 &=-120T+1-32\theta^3-32k\theta^4
                      -16k^2\theta^5-\frac{16}{3}k^3\theta^6.
 \end{split}                                                  \tag{N27.3}
\]
Here $\theta T=0$ by HD-N23. Thus these actual restriction
maps have the required nonzero Weil component at the
specified Gaussian determinant $[-7]$. The trace calculation
below applies to every choice of their nonzero sections,
including sections that vanish at intersection points.

\paragraph{1. Cohomology and the local Hom maps.}
Every $C_\nu$ is integral and $s_\nu$ is nonzero, so
multiplication by $s_\nu$ is injective on
$\mathcal O_{C_\nu}$. Therefore
\[
           \mathcal H^0(F)=\mathcal I_Z,\qquad
           \mathcal H^1(F)=\operatorname{coker}(q).       \tag{N27.4}
\]
The second sheaf is torsion; its support consists of the
zero divisors of the sections and possibly the nodes.
The first is torsion-free of rank one. The actual canonical
two-cohomology triangle then gives
$\mathcal E xt^{-1}(F,F)=0$: its possible term is
$\mathcal H om(\mathcal H^1(F),\mathcal H^0(F))=0$.
All lower self-Ext sheaves vanish by the cohomological
amplitude. The same argument, or the local-to-global
sequence, gives $\operatorname{Ext}^{<0}(F,F)=0$.
Perfection follows from smoothness of the ambient variety.

Write
\[
 A_0=\bigoplus_\nu j_{\nu,*}\mathcal O_{C_\nu},\qquad
 A_1=\bigoplus_\nu j_{\nu,*}N_{C_\nu/A},\qquad
 T_s=\bigoplus_\nu j_{\nu,*}
                         (M_\nu/s_\nu\mathcal O_{C_\nu}).
                                                               \tag{N27.5}
\]
The regular-embedding Koszul resolutions, with local
trivializations of the lines $M_\nu$, give
\[
 \begin{gathered}
 \mathcal H om(Q,Q)=A_0,\qquad
 \mathcal E xt^1(Q,Q)=A_1,\\
 \mathcal E xt^j(Q,\mathcal O_A)=0\quad(j\ne3),
 \qquad Q^\dagger:=\mathcal E xt^3(Q,\mathcal O_A)
   =\bigoplus_\nu j_{\nu,*}
                         (M_\nu^{-1}\otimes\det N_{C_\nu/A}).
 \end{gathered}                                               \tag{N27.6}
\]
There are no cross terms in the first two identities.
At a transverse node the two modules are
$R/(x_1,x_2,x_3)$ and $R/(y_1,y_2,y_3)$ in the smooth
six-variable ring; their cross Ext is concentrated in
degree three, by the actual Koszul calculation of HD-N25.
At a smooth component the usual regular-embedding Koszul
Hom complex gives its normal bundle in degree one.
Changing a line trivialization multiplies both arguments
and cancels on self-Hom, so these identifications glue
canonically. The dual calculation in (N27.6) uses the
top exterior normal line and is also canonical.

Apply sheaf derived Hom from (N27.2) into $\mathcal O_A$.
It gives
\[
 \mathcal H om(F,\mathcal O_A)=\mathcal O_A,
 \qquad \mathcal E xt^1(F,\mathcal O_A)=0,
 \qquad \mathcal E xt^2(F,\mathcal O_A)=Q^\dagger.
                                                               \tag{N27.7}
\]
The first identification sends a function to that function
times $v$. Applying Hom into $Q$, the map
$\mathcal H om(Q,Q)\to\mathcal H om(\mathcal O_A,Q)$
is the literal componentwise multiplication
\[
                  A_0\longrightarrow Q,
                    (a_\nu)_\nu\longmapsto(a_\nu s_\nu)_\nu.
                                                               \tag{N27.8}
\]
It is injective and has cokernel $T_s$. Its next Hom
sequence therefore gives
\[
 \mathcal H om(F,Q[-1])=0,\qquad
 0\longrightarrow T_s\longrightarrow\mathcal H om(F,Q)
                       \longrightarrow A_1\longrightarrow0.
                                                               \tag{N27.9}
\]
Finally apply Hom from $F$ to (N27.2). The map
$\mathcal H om(F,\mathcal O_A)\to\mathcal H om(F,Q)$
is zero, since every local map in its source is a function
times $v$ and $qv=0$. Equations (N27.7)--(N27.9) give
the exact, natural identities
\[
 \mathcal H om(F,F)=\mathcal O_A,
 \qquad
 0\longrightarrow T_s\longrightarrow\mathcal E xt^1(F,F)
                          \longrightarrow A_1\longrightarrow0.
                                                               \tag{N27.10}
\]
This identifies the actual maps through the triangle;
no splitting of this extension is asserted.
Because the negative self-Ext sheaves vanish,
$\operatorname{End}(F)=H^0(\mathcal O_A)=\mathbf C$.
The global version of (N27.6), in degrees zero and one,
also proves
\[
                      \operatorname{Hom}(F,\mathcal O_A)
                                             =\mathbf C v.
                                                               \tag{N27.11}
\]

\paragraph{2. The cohomology estimate on any such quotient arrangement.}
The same argument applies to a nodal union of $N'$ translated
abelian threefolds in any abelian sixfold, with one ample
line and nonzero section on each component. Its normal
bundles are trivial of rank three. Hence
\[
                         h^1(A_1)=9N'.                    \tag{N27.12}
\]
For an ample line $M$ on an abelian threefold $C$,
$H^i(C,M)=0$ for $i>0$. The exact primary locator is
Kodaira, \emph{On a differential-geometric method in the
theory of analytic stacks}, PNAS 39 (1953), 1268--1273,
Theorem 3 on printed p.1272. Its hypothesis is satisfied:
a very ample power of $M$ supplies a positive representative
of $c_1(M)$ by the Fubini--Study form divided by that power,
and $K_C$ is trivial. Translation-invariant Dolbeault forms
give $h^3(\mathcal O_C)=1$ and $h^1(\mathcal O_C)=3$.
Thus the actual sequence
$0\to\mathcal O_C\xrightarrow{s}M\to M/s\mathcal O_C\to0$
gives
\[
                H^2(C,M/s\mathcal O_C)
                     \simeq H^3(C,\mathcal O_C),
                   \qquad h^2(T_s)=N'.                   \tag{N27.13}
\]
This does not require the zero divisor to be smooth.
For the existence assertion above, the same vanishing and
abelian Riemann--Roch give
$h^0(M)=c_1(M)^3/6>0$.

The cohomology sequence of (N27.10) now implies
\[
                 h^1(\mathcal E xt^1(F,F))\ge8N'.         \tag{N27.14}
\]
In the actual local-to-global self-Ext sequence, the term
$E_2^{1,1}$ has no incoming differential. Its only possible
outgoing differential is $d_2$ into
$H^3(\mathcal H om(F,F))=H^3(\mathcal O)$, of dimension 20;
all later targets have negative self-Ext sheaf degree.
Consequently
\[
                        \dim\operatorname{Ext}^2(F,F)
                                             \ge8N'-20.
                                                               \tag{N27.15}
\]
The full degree-two trace target on an abelian sixfold has
dimension 495. Its ordinary trace kernel therefore has
dimension at least $8N'-515$.

\paragraph{3. Every coherent action is compared with an actual descended pair.}
First work on the original $A$, or on any isogeny pullback
$f:X\to A$ with its pulled arrangement and data. Continue
to write $F,Z,Q,v$ for the resulting objects on $X$.
For any point functor $T_a^*(-)\otimes P$ preserving $F$,
(N27.4) gives an isomorphism
$T_a^*\mathcal I_Z\otimes P\simeq\mathcal I_Z$.
Taking double duals gives $P\simeq\mathcal O_X$; extending
the ideal isomorphism to this line proves $T_a(Z)=Z$.
Thus a faithful finite point subgroup is an actual finite
translation subgroup $G$ preserving $Z$. For an isogeny
pullback only the inclusion
\[
                            f(G)\subseteq\operatorname{Stab}(Z_{\rm original})
                                                               \tag{N27.16}
\]
is used. No equality with a stabilizer of the original
complex is asserted.

Choose the trivializations of the lines $P$ just obtained,
and write $u_a:T_a^*F\to F$ for the given isomorphisms.
By (N27.11) there is a unique $\lambda_a\in\mathbf C^*$
with
\[
                  v u_a=\lambda_a T_a^*v,
                    \qquad u'_a=\lambda_a^{-1}u_a.       \tag{N27.17}
\]
Both $u'_{a+b}$ and $u'_aT_a^*u'_b$ are isomorphisms
$T_{a+b}^*F\to F$ whose composition with $v$ is
$T_{a+b}^*v$. This Hom space has dimension one by
simplicity, and composition with $v$ is nonzero. They
are therefore equal. The $u'_a$ give the literal
translation cocycle, with $v$ equivariant for the standard
action on $\mathcal O_X$.
More explicitly, if the original functor composition
has scalar cocycle $c_{a,b}$, composition with $v$ shows
$\lambda_a\lambda_b=c_{a,b}\lambda_{a+b}$, in the
convention $u_aT_a^*u_b=c_{a,b}u_{a+b}$.
Thus (N27.17) supplies the actual scalar cochain comparison.
Conjugation on Ext is unchanged. Gluability makes all
negative self-Ext vanish, so there is no additional higher
automorphism coherence to choose.

We spell out a canonical sheaf reconstruction of $Q$ and
$q$ from this equivariant pair $(F,v)$. Put
$R=\mathcal E xt^2(F,\mathcal O_X)=Q^\dagger$.
The codimension-three Koszul duality in (N27.6) gives
\[
             R\mathcal H om(R,\mathcal O_X)=Q[-3],
                    \qquad Q=\mathcal E xt^3(R,\mathcal O_X).
                                                               \tag{N27.18}
\]
These are the canonical double-dual maps: on a component,
the two top normal exterior lines cancel and the two
line duals return $M$; they glue by the Koszul evaluation
map. The dual triangle of (N27.2) is
\[
             R[-3]\xrightarrow{q^\vee}\mathcal O_X
                  \xrightarrow{v^\vee}F^\vee
                  \longrightarrow R[-2].                \tag{N27.19}
\]
Its rotation is the canonical two-cohomology triangle of
$F^\vee$, with $\mathcal H^0(F^\vee)=\mathcal O_X$
identified by $v^\vee$ and $\mathcal H^2(F^\vee)=R$.
Write its connecting map as
$\delta:R[-2]\to\mathcal O_X[1]$. Rotation gives
$\delta=-q^\vee[1]$, and the canonical double-dual map gives
$q=-(\delta[-1])^\vee:\mathcal O_X\to Q$.
Thus every $u'_a$
induces an honest, coherently compatible action on the
sheaf $Q$ and its section $q$. This proves equivariance
of the presented two-term complex without choosing lifts
on a cone or assuming the original sections descend.

Let $p:X\to Y=X/G$ be the translation quotient isogeny.
The literal equivariant sheaf and section descend by
Stacks, Proposition 35.5.2 (Tag 023T) and its finite-etale
form Lemma 35.6.2 (Tag 0D1V). Explicitly
$Q'=(p_*Q)^G$, and the section is the descended map
$\mathcal O_Y\to Q'$. The identity
$X\times_YX\simeq G\times X$ identifies its descent
cocycle with (N27.17)--(N27.19), and the descent unit gives
$p^*Q'\simeq Q$ with the original section.

The action on $Q$ permutes its actual support components;
there are no Hom maps between distinct supports.
For one component $C$ and its stabilizer $G_C$, the
quotient $C\to C/G_C$ is a translation isogeny. Its
equivariant line descends to an actual line $M'$ on
$C/G_C$, and the section descends to a nonzero section.
The image of $C/G_C$ is an embedded translated abelian
threefold in $Y$: its inverse image is the union of the
parallel, pairwise disjoint translates of $C$. Thus
\[
                     Q'=\bigoplus_{\text{component orbits}}
                                  j'_*M'.                 \tag{N27.20}
\]
Each $M'$ is ample, because its finite surjective pullback
is the given ample $M$; the exact source is Stacks,
Lemma 30.17.2 (Tag 0B5V). The quotient is etale, so distinct
components have exactly the same transverse double-point
local geometry and no triple point. We obtain the actual
descended object
\[
                    F'=[\mathcal O_Y\longrightarrow Q'],
                               \qquad p^*F'\simeq F.      \tag{N27.21}
\]

\paragraph{4. Component counts and the invariant trace kernel.}
On the original arrangement the translation stabilizer has
order at most 144 by the actual two-direction calculation
of HD-N23. Hence every quotient by a preserving subgroup
has at least $\lceil9408/144\rceil=66$ component orbits.
After an arbitrary isogeny pullback, the map from connected
inverse-image components to their original labels is
surjective and is equivariant for (N27.16). It induces a
surjection on orbit sets. Thus the actual quotient
arrangement in (N27.20) still has $N'\ge66$ components,
even when the new translation subgroup has larger order.
Only the image group in the original support stabilizer
enters this estimate.

Apply (N27.15) to $F'$ and its descended ample lines.
It has ordinary degree-two trace kernel of dimension at least
\[
                         8\cdot66-515=13.                 \tag{N27.22}
\]
Finite-etale pullback identifies its self-Ext with the
canonical invariant self-Ext of $F$: the adjunction and
projection formula give
\[
 \operatorname{Ext}^k_Y(F',F')
           \xrightarrow{\ \sim\ }
                    \operatorname{Ext}^k_X(p^*F',p^*F')^G,
                                                               \tag{N27.23}
\]
since $(p_*\mathcal O_X)^G=\mathcal O_Y$ and averaging
is exact. Pullback of the Atiyah class and its trace is
natural for the etale map; cohomological pullback is
injective, split by the trace divided by $|G|$.
Thus all thirteen kernel dimensions remain in the actual
invariant trace kernel upstairs. Rescaling (N27.17) left
this conjugation action unchanged, so the statement applies
to the originally given coherent action as well.
The identity-sector Hochschild comparison of HD-N12
identifies the same weak categorical trace kernel.

Subsequent actual line tensor and Fourier--Mukai equivalences
carry the point actions and trace maps by the intertwiners
of HD-N12 and HD-N17. The scope is the actual objects
(N27.1), their arbitrary isogeny pullbacks and these
subsequent equivalences. No equality of object stabilizers
under isogeny, no arbitrary lift of a projector, and no
descent of a section without the displayed maps is used.
